\section{Frontend and Abstract Syntax Tree Construction}\label{sec:implementation-frontend}

The frontend validates syntactic structure and builds the logical model of the composition.
It is implemented as a classic Flex/Bison scanner--parser pipeline that builds an \texttt{ast::Program}
consumed by later
stages.

\subsection{Entry Points and Result Type}\label{detail-subsec:frontend-entry-points}

The parsing stage exposes two public entry points:

\topic{\texttt{parse\_stream(FILE*, source\_name, diagnostics)}.} Used by \texttt{motivo::compile}; wraps
the input file in a \texttt{ScannerInputGuard} that resets Flex state, attaches the shared
\texttt{DiagnosticsEngine}, and calls \texttt{yyrestart}.

\topic{\texttt{parse\_source(string, source\_name, diagnostics)}.} Used in unit tests; scans an in-memory
string via \texttt{yy\_scan\_string}.

Both paths construct an empty \texttt{ast::Program}, invoke \texttt{detail::Parser(loc, diagnostics,
program).parse()}, and return a \texttt{ParseResult} owning the completed tree.
If Bison returns \texttt{EXIT\_FAILURE} without prior diagnostics, the parser emits a generic parse error.

\subsection{Lexical Analysis: Identifying Musical
Literals}\label{detail-subsec:lexical-analysis:-identifying-musical-literals}

The scanner is defined in \texttt{parsing/detail/scanner.l} and generated by Flex at build time.
It transforms the character stream into Bison token values (\texttt{Parser::symbol\_type}) while maintaining a
\texttt{source::Location} per token via \texttt{YY\_USER\_ACTION} column/line updates.

\topic{Pitch-Class and Note Recognition.}

Unlike conventional languages that treat musical data as strings, Motivo recognizes pitch literals such as
\texttt{F\#4} as primitive constants.
The scanner uses dedicated character-class macros:

\begin{lstlisting}[caption={Flex regular expression macros for note literal recognition (\texttt{scanner.l}).},label={detail-lst:scanner-note-macros}]
pitch        [A-G]
accidental   [#b]
octave       [0-8]
note_lit     {pitch}{accidental}?{octave}
\end{lstlisting}

When \texttt{note\_lit} matches, the action calls \texttt{parse\_note\_literal} from
\texttt{parsing/detail/scanner\_helpers.cpp}, which constructs a \texttt{music::Note} (pitch enum, accidental enum,
octave integer) and returns a \texttt{TOK\_NOTE\_LIT} token carrying that value.
Drum literals (\texttt{kick}, \texttt{snare}, etc.) and instrument keywords (\texttt{piano}, \texttt{drums}) are
recognized in the same file and mapped to \texttt{music::DrumNote} and \texttt{music::Instrument} enums.

\topic{Keyword and Symbol Disambiguation.}

Reserved keywords include structural tokens (\texttt{track}, \texttt{pattern}, \texttt{voice}, \texttt{play}) and
type keywords (\texttt{int}, \texttt{double}, \texttt{bool}, \texttt{note}, \texttt{seq}, \texttt{chord}) used in
\texttt{VarDeclStatement} and \texttt{TypedParameter} positions.
Identifiers match \texttt{\{alpha\}\{alnum\}*}; because keywords are listed before the identifier rule, longest-match
disambiguation ensures names such as \texttt{loop\_count} are not tokenized as the \texttt{loop} keyword.

Lexical errors call \texttt{report\_lexical}, which forwards to \texttt{DiagnosticsEngine::report} with
\texttt{DiagnosticStage::Parsing} when a diagnostics pointer is set by \texttt{scanner\_set\_diagnostics}.

\subsection{Syntactic Analysis: Building the Hierarchy}\label{detail-subsec:syntactic-analysis:-building-the-hierarchy}

The parser grammar lives in \texttt{parsing/detail/parser.y} and is processed by Bison~3.8 with C++ variant-based
semantic values (\texttt{\%define api.value.type variant}).
The generated \texttt{motivo::parsing::detail::Parser} class receives three parse parameters: a
\texttt{source::Location}, a \texttt{DiagnosticsEngine\&}, and the \texttt{ast::Program\&} under construction.

Grammar actions use helper macros \texttt{expression(...)} and \texttt{statement(...)} that allocate
\texttt{std::unique\_ptr\textless ast::Expression\textgreater{}} and \texttt{std::unique\_ptr\textless
ast::Statement\textgreater{}} nodes.
The start symbol \texttt{program} assembles \texttt{program.header}, \texttt{program.globals}, and
\texttt{program.tracks} from nested \texttt{track\_def}, \texttt{pattern\_def}, and \texttt{voice\_def} productions.
Detailed productions mirror the EBNF documented in \texttt{compiler/docs/grammar.bnf}.

\topic{AST Node Architecture: Classes and Sum-Types.}

The AST lives in the \texttt{motivo::ast} namespace.
Nodes are grouped into four structural classes rather than catalogued individually: \emph{program shell},
\emph{definitions}, \emph{statements}, and \emph{expressions}.
Within each class, \texttt{std::variant} (\texttt{StatementKind}, \texttt{ExpressionKind}, \texttt{TrackItem},
\texttt{GlobalItem}) provides closed sum types; \texttt{std::visit} in later stages dispatches without virtual methods.

\topic{Program shell.}
\texttt{ast::Program} is the root owned by \texttt{ParseResult}.
It contains three fields: a \texttt{Header} (optional \texttt{tempo} and \texttt{signature} declarations), a
\texttt{globals} vector of \texttt{GlobalItem} nodes, and a \texttt{tracks} vector of \texttt{TrackDefinition} nodes.
The header is validated semantically before any track body is visited.

\topic{Definition nodes.}
Definitions establish lexical boundaries and appear as top-level or nested items:

\topic{\texttt{TrackDefinition}.} Optional track name, optional \texttt{music::Instrument} from a
\texttt{using} clause, and a \texttt{body} of \texttt{TrackItem} variants (statement, local pattern, or voice).

\topic{\texttt{PatternDefinition}.} Name, typed parameter list (\texttt{TypedParameter} pairs a
\texttt{types::Type} with an identifier), and a statement block.
Patterns may appear at global, track, or voice scope; overload discrimination uses the full parameter-type signature.

\topic{\texttt{VoiceDefinition}.} Optional \texttt{from} offset expression and a body of \texttt{VoiceItem}
nodes (statements or local patterns).
Voices introduce parallel temporal lanes without duplicating the AST node taxonomy.

\topic{Statement nodes.}
All executable actions share a \texttt{Statement} wrapper carrying \texttt{StatementKind} and \texttt{source::Location}.
The variant distinguishes three families:

\topic{Declarative.} \texttt{VarDeclStatement} (typed name and initializer) and \texttt{AssignStatement}
(name and value expression).
Version~2 replaces v1's untyped \texttt{LetStatement} with explicit \texttt{types::Type} on declarations.

\topic{Control flow.} \texttt{ForStatement} (optional init, condition, step, body), \texttt{LoopStatement}
(integer count and body), and \texttt{IfStatement} (condition, then-branch, optional else-branch).
These nodes preserve the source control structure until lowering unrolls them.

\topic{Musical.} \texttt{PlayStatement} wraps a \texttt{PlayTarget}: a \texttt{PlaySource} variant
(expression or drum literal), optional duration expression (\texttt{:} operand), and optional \texttt{from} offset.
The target also carries its own location for diagnostics on duration and offset operands.

\topic{Expression nodes.}
Expressions are the most numerous node family; the thesis focuses on the categories lowering and semantic analysis
dispatch on, not every literal wrapper:

\topic{Scalar literals.} \texttt{IntLiteralExpression}, \texttt{FloatLiteralExpression},
\texttt{BoolLiteralExpression}.

\topic{Musical literals and composites.} \texttt{NoteLiteralExpression}, \texttt{RestLiteralExpression},
\texttt{DrumNoteLiteralExpression}; \texttt{SequenceExpression} and \texttt{ChordExpression} as vectors of
\texttt{DurationalTarget} pairs (value expression plus optional duration).

\topic{Operators.} \texttt{UnaryExpression}, \texttt{BinaryExpression}, \texttt{TernaryExpression}, and
\texttt{ParenthesisedExpression} (delegates to an inner expression).

\topic{Name and call resolution.} \texttt{IdentifierExpression} and \texttt{PatternCallExpression} (callee
name plus argument expression list).
Semantic analysis attaches resolved \texttt{SymbolId}s and types to these nodes in \texttt{Annotations}.

\topic{Ownership and locations.}
Recursive subtrees use \texttt{std::unique\_ptr} (\texttt{StatementPtr}, \texttt{ExpressionPtr}) so the AST remains a
tree in memory.
Every \texttt{Statement} and \texttt{Expression} carries a \texttt{source::Location} copied from the Bison \texttt{loc}
parameter at reduction time, which drives diagnostic messages in semantic analysis and lowering failures.

The combination of class-level grouping, variant dispatch, and pointer-based recursion provides the foundation for the
single-pass semantic traversal and the interpreter-style lowering pass that follow.
